\subsection{拓扑空间的附着}

设 X、B 为拓扑空间，设 A 为 B 的子空间，并设 \(f: A \rightarrow X\) 为连续映射。用 Y 表示空间 \(Y_{1} = X\) 与 \(Y_{2} = B\) 的拓扑和，并借助典范单射把 \(Y_{1}\) 与 \(Y_{2}\) 等同于 Y 的子空间。设 R 为 Y 上最细的等价关系，使得 \(Y_{2}\) 的元素 a 与 \(f(a)\)（即 \(Y_{1}\) 的元素）等价（对每个 \(a \in A\)）。关系 R 在 \(Y_{1}\) 上就是相等关系。设 \(x \in Y_{1}\) 与 \(b \in Y_{2}\)；我们有 x R b 当且仅当 \(b \in A\) 且 \(f(b) = x\)。设 \(b, b' \in Y_{2}\)；要使 b R \(b'\) 成立，必须而且只需 \(b = b'\)，或者 \(b, b' \in A\) 且 \(f(b) = f(b')\)。

定义 9. --- 商拓扑空间 Y/R 称为沿映射 f 把空间 B 附着到空间 X 上所得的空间，记作 \(X \cup_{f} B\)。

用 \(\alpha_{f}\colon\mathrm{X}\to\mathrm{X}\cup_{f}\mathrm{B}\) 与 \(\beta_{f}\colon\mathrm{B}\to\mathrm{X}\cup_{f}\mathrm{B}\) 表示 Y 到 Y/R 上的典范满射在 \(\mathrm{Y}_{1}\) 与 \(\mathrm{Y}_{2}\) 上的限制。我们有 \(\alpha_{f}\circ f=\beta_{f}\mid\mathrm{A}\)。

注记。--- 1) 映射 \(\alpha_{f}\) 是单射。对 X 的每个子集 U，有 \(\bar{\alpha}_{f}^{-1}(\alpha_{f}(U)) = U\) 与 \(\bar{\beta}_{f}^{-1}(\alpha_{f}(U)) = f^{-1}(U)\)。若 A 是 B 的开（相应地，闭）子集，这些关系式蕴涵 \(\alpha_{f}\) 是从 X 到 \(X \cup_{f} B\) 中的开（相应地，闭）映射。

设 U 是 X 的开子集，V 是 B 的开集，且满足 \(f^{-1}(U) = V \cap A\)。\(Y_{1}\) 的子集 U 与 \(Y_{2}\) 的子集 V 的并是 Y 的开饱和子集；因此它在 \(X \cup_{f} B\) 中的像是开集，它在 \(\alpha_{f}(X)\) 上的迹是 \(\alpha_{f}(U)\)。于是，映射 \(\alpha_{f}\) 定义了 X 到它在 \(X \cup_{f} B\) 中的像上的一个同胚。

\begin{enumerate}
\def\labelenumi{\arabic{enumi})}
\setcounter{enumi}{1}
\tightlist
\item
  设 V 是 B 的一个子集。我们有 \(\bar{\alpha}_{f}^{-1}(\beta_{f}(V)) = f(V \cap A)\) 与 \(\bar{\beta}_{f}^{-1}(\beta_{f}(V)) = V \cup \bar{f}(f(V \cap A))\)。若 A 在 B 中闭且映射 f 是闭映射，则映射 \(\beta_{f}\) 是闭映射。类似地，若 A 在 B 中开且映射 f 是开映射，则映射 \(\beta_{f}\) 是开映射。
\end{enumerate}

在这两种情形下，映射 \(\beta_{f}\) 都诱导出 \(B\setminus A\) 到其像上的一个同胚。

\begin{enumerate}
\def\labelenumi{\arabic{enumi})}
\setcounter{enumi}{2}
\tightlist
\item
  假设 A 是闭的且映射 \(f\) 是逆紧的。我们刚才已看到映射 \(\beta_f\) 是闭映射。对每个 \(x \in X\)，有 \(\overline{\beta_f}(\alpha_f(x)) = \overline{f}(x)\)。因此它是 A 的一个拟紧子集（TG, I, p. 75, 定理 1），从而也是 B 的拟紧子集。对每个 \(b \in B\)，\(\overline{\beta_f}(\beta_f(b))\) 等于 \(\{b\}\)（若 \(b \in B - A\)），或等于 \(\overline{f}(f(b))\)（若 \(b \in A\)）；两种情形下它都是 B 的拟紧子集。因此映射 \(\beta_{f}\) 的纤维都是拟紧的；从而（同前所引）映射 \(\beta_{f}\) 是逆紧的。
\end{enumerate}

映射 \(\alpha_{f}\) 也是逆紧的（TG, I, p. 72, 命题 2）。由此推出，从 Y 到 X \(\cup_{f}\) B 上的典范映射是逆紧的。

例。--- 1) 设 X 与 Y 为拓扑空间，\(f\colon X\to Y\) 为连续映射。柱 \(\operatorname{Cyl}(f)\) 无非就是把空间 \(X\times I\) 沿映射 \(f\circ \mathrm{pr}_1\)（从 \(X\times \{1\}\) 到 Y 中）附着到空间 Y 上所得的空间。

\begin{enumerate}
\def\labelenumi{\arabic{enumi})}
\setcounter{enumi}{1}
\tightlist
\item
  设 X 为拓扑空间，B 为拓扑空间，A 为 B 的子空间，并设 \(f: A \to X\) 为连续映射，它诱导出 A 到其像上的一个同胚。
\end{enumerate}

令 \(\mathrm{A}' = f(\mathrm{A})\)；设 \(f'\) 是从 \(\mathrm{A}'\) 到 B 中的映射，它把每个 \(x \in \mathrm{A}'\) 对应到 \(x\) 在 \(f\) 下的唯一原像。映射 \(f' \circ f\) 是连续的；由于 \(f\) 是严格映射，映射 \(f'\) 是连续的。存在唯一的连续映射 \(u\)，从空间 \(\mathrm{X} \cup_f \mathrm{B}\) 到空间 \(\mathrm{B} \cup_{f'} \mathrm{X}\)，它把 \(\alpha_f(x)\) 映为 \(\beta_{f'}(x)\)（\(x \in \mathrm{X}\)），把 \(\beta_f(b)\) 映为 \(\alpha_{f'}(b)\)（\(b \in \mathrm{B}\)）；它是一个同胚，把子空间 \(\alpha_f(\mathrm{X})\) 映到子空间 \(\beta_{f'}(\mathrm{X})\) 上。

命题 8. --- 设 Z 为拓扑空间。

\begin{enumerate}
\def\labelenumi{\alph{enumi})}
\item
  设 \(u \colon X \to Z\) 与 \(v \colon B \to Z\) 为连续映射，满足 \(v(a) = u(f(a))\)（对所有 \(a \in A\)）。于是存在唯一的连续映射 \(w \colon X \cup_f B \to Z\)，使得 \(w \circ \alpha_f = u\) 且 \(w \circ \beta_f = v\)。
\item
  设 \(\sigma\colon\mathbf{X}\times\mathbf{I}\to\mathbf{Z}\) 与 \(\tau\colon\mathbf{B}\times\mathbf{I}\to\mathbf{Z}\) 为同伦。假设 \(\tau\) 在 A 上是固定的，且 \(\sigma(f(a),t)=\tau(a,t)\)（对所有 \(a\in\mathbf{A}\) 及每个 \(t\in\mathbf{I}\)）。于是存在唯一的同伦 \(\eta\colon(\mathbf{X}\cup_{f}\mathbf{B})\times\mathbf{I}\to\mathbf{Z}\)，使得 \(\eta(\alpha_{f}(x),t)=\sigma(x,t)\) 与 \(\eta(\beta_{f}(b),t)=\tau(b,t)\)（对 \(x\in\mathbf{X}\)、\(b\in\mathbf{B}\) 与 \(t\in\mathbf{I}\)）。这个同伦在 \(\beta_{f}(\mathbf{A})\) 上是固定的。
\end{enumerate}

本命题由商空间的定义及命题 10（I, p. 19）立即得出。

推论 1. --- 设 \(B'\) 为拓扑空间，\(A'\) 为 \(B'\) 的子空间，\(f': A' \to X\) 为连续映射。设 \(v: B \to B'\) 为连续映射，满足 \(v(A) \subset A'\) 且 \(f = f' \circ (v \mid A)\)。设 \(w: X \cup_f B \to X \cup_{f'} B'\) 为唯一的连续映射，满足 \(w \circ \alpha_f = \alpha_{f'}\) 且 \(w \circ \beta_f = \beta_{f'} \circ v\)。若 \(v\) 定义了偶 \((B, A)\) 到偶 \((B', A')\) 上的同伦等价，则 \(w\) 定义了偶 \((X \cup_f B, \alpha_f(X))\) 到偶 \((X \cup_{f'} B', \alpha_{f'}(X))\) 上的同伦等价。

设 \(v' \colon \mathrm{B}' \to \mathrm{B}\) 为连续映射，\(\tau \colon \mathrm{B} \times \mathbf{I} \to \mathrm{B}\) 为在 A 上固定、连接 \(\mathrm{Id}_{\mathrm{B}}\) 与 \(v' \circ v\) 的同伦，\(\tau' \colon \mathrm{B}' \times \mathbf{I} \to \mathrm{B}'\) 为在 \(\mathrm{A}'\) 上固定、连接 \(\mathrm{Id}_{\mathrm{B}'}\) 与 \(v \circ v'\) 的同伦。对 \(a' \in \mathrm{A}'\) 与 \(a = v'(a')\)，我们有 \(a' = v(a)\) 且 \(\beta_f(v'(a')) = \beta_f(a) = \alpha_f(f(a)) = \alpha_f(f'(a'))\)。由命题 8(a)，存在唯一的映射 \(w' \colon X \cup_{f'} B' \to X \cup_f B\)，满足 \(w' \circ \alpha_{f'} = \alpha_{f'}\) 且 \(w' \circ \beta_{f'} = \beta_f \circ v'\)。由命题 8(b)，存在唯一的同伦 \(\eta \colon (\mathrm{X} \cup_f \mathrm{B}) \times \mathbf{I} \to X \cup_{f'} B'\)，满足 \(\eta(\alpha_f(x), t) = \alpha_{f'}(x)\) 与 \(\eta(\beta_f(b), t) = \beta_{f'}(\tau(b, t))\)（对 \(x \in X\)、\(b \in B\) 与 \(t \in I\)）。同样存在唯一的同伦 \(\eta' \colon (X \cup_{f'} B') \times I \to X \cup_f B\)，满足 \(\eta'(\alpha_{f'}(x), t) = \alpha_f(x)\) 与 \(\eta'(\beta_{f'}(b), t) = \beta_f(\tau'(b, t))\)（对 \(x \in X\)、\(b \in B'\) 与 \(t \in I\)）。于是，从 \((X \cup_f B) \times I\) 到 \(X \cup_f B\) 中的映射 \((x, t) \mapsto \eta'(\eta(x, t), t)\) 是在 \(\alpha_f(X)\) 上固定、连接 \(X \cup_f B\) 的恒同映射与映射 \(w' \circ w\) 的同伦。类似地，从 \((X \cup_{f'} B') \times I\) 到 \(X \cup_{f'} B'\) 中的映射 \((x, t) \mapsto \eta(\eta'(x, t), t)\) 是在 \(\alpha_{f'}(X)\) 上固定、连接 \(X \cup_{f'} B'\) 的恒同映射与映射 \(w \circ w'\) 的同伦。推论由此得出。

例 3. --- 用 i 表示 A 到 B 中的典范单射，取映射柱 i 作为空间 B'；用 \(\alpha_{i}\colon\mathrm{A}\times\mathbf{I}\to\mathrm{Cyl}(i)\) 与 \(\beta_{i}\colon\mathrm{B}\to\mathrm{Cyl}(i)\) 表示典范映射，用 \(r_{i}\colon\mathrm{Cyl}(i)\to\mathrm{B}\) 表示柱 \(\mathrm{Cyl}(i)\) 到其柱底上的典范收缩映射。设 \(A_{0}\) 为 \(\alpha_{i}(\mathrm{A}\times\{0\})\)，即 \(\mathrm{Cyl}(i)\) 的子空间，并用 \(f_{0}\colon\mathrm{A}_{0}\to\mathrm{X}\) 表示映射 \(f\circ r_{i}\)。对 \(a\in A\)，我们有

\[
\beta_ {f} \circ r _ {i} (\alpha_ {i} (a, 0)) = \beta_ {f} (a) = \alpha_ {f} (f _ {0} (\alpha_ {i} (a, 0))) = \alpha_ {f} (f (a)).
\]

设 \(\eta\colon X\cup_{f_0}\mathrm{Cyl}(i)\to X\cup_fB\) 为唯一的连续映射，满足 \(\eta\circ\alpha_{f_0}=\alpha_f\) 且 \(\eta\circ\beta_{f_0}=\beta_f\circ r_i\)。假设偶 (B, A) 具有同伦扩张性质。于是映射 \(r_i\) 定义了偶 \((\mathrm{Cyl}(i),\mathrm{A}_0)\) 到偶 (B,A) 上的同伦等价（III, p. 243, 定理 1）。由推论 1，映射 \(\eta\) 便定义了偶 \((X\cup_{f_0}\mathrm{Cyl}(i),\alpha_{f_0}(X))\) 到偶 \((X\cup_fB,\alpha_f(X))\) 上的同伦等价。特别地，\(\eta\) 是一个同伦等价。

推论 2. --- 设 \(X'\) 为拓扑空间，\(u: X \to X'\) 为连续映射，并令 \(f' = u \circ f\)。设 \(w: X \cup_f B \to X' \cup_{f'} B\) 为唯一的连续映射，满足 \(w \circ \alpha_f = \alpha_{f'} \circ u\) 且 \(w \circ \beta_f = \beta_{f'}\)。若 \(u\) 定义了偶 \((X, f(A))\) 到偶 \((X', f'(A))\) 上的同伦等价，则 \(w\) 定义了偶 \((\mathrm{X}\cup_f\mathrm{B},\beta_f(\mathrm{B}))\) 到偶 \((\mathrm{X}'\cup_{f'}\mathrm{B},\beta_{f'}(\mathrm{B}))\) 上的同伦等价。

证明与推论 1 的证明类似。

命题 9. --- 设 X 与 B 为拓扑空间，A 为 B 的子空间，f: A → X 为连续映射。

\begin{enumerate}
\def\labelenumi{\alph{enumi})}
\item
  若偶 (B, A) 具有同伦扩张性质，则偶 \((\mathrm{X} \cup_{f} \mathrm{B}, \alpha_{f}(\mathrm{X}))\) 亦然。
\item
  假设映射 f 是单射且严格的。若偶 \((\mathrm{X}, f(\mathrm{A}))\) 具有同伦扩张性质，则偶 \((\mathrm{X} \cup_{f} \mathrm{B}, \beta_{f}(\mathrm{B}))\) 亦然。
\item
  假设偶 (B, A) 具有同伦扩张性质。设 Z 为拓扑空间，\(u: X \cup_{f} B \to Z\) 为连续映射，\(\sigma: \alpha_{f}(X) \times I \to Z\) 为一同伦，其末项映射是 u 在子空间 \(\alpha_{f}(X)\) 上的限制。
\end{enumerate}

令 \(v_{1} = u \circ \alpha_{f}\)，并用 \(\tau_{1} \colon \mathbf{X} \times \mathbf{I} \to \mathbf{Z}\) 表示由 \((x,t) \mapsto \sigma(\alpha_{f}(x),t)\) 给出的映射。令 \(v_{2} = u \circ \beta_{f}\)。由于 \(\beta_{f} \mid \mathrm{A} = \alpha_{f} \circ f\)，从 \(\mathrm{A} \times \mathbf{I}\) 到 \(\mathrm{Z}\) 中、把 \((a,t)\) 映到 \(\sigma(\alpha_{f}(f(a)),t)\)（\(a \in \mathrm{A}\)，\(t \in \mathbf{I}\)）的映射是一个同伦，其末项等于映射 \(v_{2} \mid \mathrm{A}\)。由于偶 \((\mathrm{B},\mathrm{A})\) 具有同伦扩张性质，存在同伦 \(\tau_{2} \colon \mathbf{B} \times \mathbf{I} \to \mathbf{Z}\)，其末项为映射 \(v_{2}\)，且满足 \(\tau_{2}(a,t) = \sigma(\alpha_{f}(f(a)),t)\)（\((a,t) \in \mathrm{A} \times \mathbf{I}\)）。

对 \(a \in \mathbf{A}\) 与 \(t \in \mathbf{I}\)，我们有 \(\tau_2(a, t) = \tau_1(f(a), t)\)。由命题 8，存在唯一的连续映射 \(\sigma: (\mathrm{X} \cup_f \mathrm{B}) \times \mathbf{I} \to \mathbf{Z}\)，满足 \(\sigma(\alpha_f(x), t) = \sigma_1(x, t)\) 与 \(\sigma(\beta_f(b), t) = \sigma_2(b, t)\)（对 \(x \in \mathrm{X}\)、\(b \in \mathrm{B}\) 与 \(t \in \mathbf{I}\)）。这是一个以 \(u\) 为末项、开拓 \(\tau\) 的同伦，由此得断言 a)。

\begin{enumerate}
\def\labelenumi{\alph{enumi})}
\setcounter{enumi}{1}
\tightlist
\item
  注意到例 2（III, p. 249），断言 b) 由断言 a) 得出。
\end{enumerate}

\subsection{缩塌子空间所得的空间}

设 X 为拓扑空间，A 为 X 的子集。考虑 X 上最细的等价关系 R，使 A 的所有元素彼此等价：X 的两个元素关于这个关系等价，当且仅当它们相等，或者它们都属于 A。

定义 10. --- X 对 R 的商空间记作 X/A，称为由 X 缩塌子集 A 所得的空间。

用 \(\rho\colon X\to X/A\) 表示典范满射。设 Y 为拓扑空间，\(f\colon X\to Y\) 为连续映射。要存在连续映射 \(g\colon X/A\to Y\) 满足 \(g\circ\rho=f\)，必须而且只需 \(f\) 在 A 上是常值的。

若集合 A 在 X 中闭（相应地，开），则映射 \(\rho\) 是闭（相应地，开）映射，并诱导出 \(X\setminus A\) 到其像上的一个同胚。若 A 是 X 的闭且拟紧的子集，则映射 \(\rho\) 是逆紧的（TG, I, p. 75, 定理 1）。

若集合 A 为空，则 \(\rho\) 是同胚。若 A 非空，则 A 缩塌为 X/A 的一点，该点称为 X/A 的基点，记作 \(s_{\mathrm{X / A}}\)。这时把空间 X/A 等同于沿从 A 到 \(\{s_{\mathrm{X / A}}\}\) 的唯一映射把空间 X 附着到空间 \(\{s_{\mathrm{X / A}}\}\) 上所得的空间。

命题 10. --- 设 X 为拓扑空间，A 为 X 的子集。假设存在同伦 \(\sigma: \mathrm{X} \times \mathbf{I} \to \mathrm{X}\)，以 \(\mathrm{Id}_{\mathrm{X}}\) 为初始，在 \(\mathrm{A} \times \{1\}\) 上是常值的，且满足 \(\sigma(\mathrm{A} \times \mathbf{I}) \subset \mathrm{A}\)。那么从 X 到 X/A 的典范映射 \(\rho\) 是同伦等价。

用 \(f\) 表示 \(\sigma\) 的末项。它是从 X 到 X 中的连续映射，同伦于 \(\mathrm{Id}_{\mathrm{X}}\)。它在 A 上是常值的；因此存在唯一的连续映射 \(g \colon \mathrm{X} / \mathrm{A} \to \mathrm{X}\)，满足 \(g \circ \rho = f\)。另一方面，由于 \(\sigma(\mathrm{A} \times \mathbf{I}) \subset \mathrm{A}\)，存在映射 \(\sigma'\)，它从 \(\mathrm{X} / \mathrm{A} \times \mathbf{I}\) 到 \(\mathrm{X} / \mathrm{A}\) 中，满足 \(\sigma'(\rho(x), t) = \rho(\sigma(x, t))\)（对所有 \(x \in \mathrm{X}\) 及每个 \(t \in \mathbf{I}\)）。由 I, p. 19, 命题 10，映射 \(\sigma'\) 是连续的。它是连接 \(\mathrm{X} / \mathrm{A}\) 的恒同映射与映射 \(\rho \circ g\) 的同伦。因此，映射 \(g\) 与 \(\rho\) 是互为同伦逆的同伦等价。

注记。--- 设 X 为拓扑空间，A 为 X 的子集。

\begin{enumerate}
\def\labelenumi{\arabic{enumi})}
\item
  若从 X 到 X/A 上的典范映射是同伦等价，则存在同伦 \(\sigma\colon\mathrm{X}\times\mathbf{I}\to\mathrm{X}\)，它以 \(\mathrm{Id}_{\mathrm{X}}\) 为初始，且在 \(\mathrm{A}\times\{1\}\) 上为常值。但也可能不存在除此之外还满足条件 \(\sigma(\mathrm{A}\times\mathbf{I})\subset\mathrm{A}\) 的同伦（III, p. 322, 习题 4）。
\item
  可能存在连接 X 的恒同映射与 X 到 X 中、在 A 上为常值的映射的同伦，而空间 X 与 X/A 却不同伦等价（III, p. 325, 习题 14）。
\item
  假设 A 是可缩的。若偶 (X, A) 具有同伦扩张性质，则从 X 到 X/A 的典范映射是同伦等价。事实上，设 \(\sigma\colon A \times I \to A\) 为一同伦，其初始映射是 A 的恒同映射，末项映射是常值映射。于是存在同伦 \(\sigma'\)，它以 \(\mathrm{Id}_{X}\) 为初始，且是 \(\sigma\) 的开拓。断言便由命题 10 得出。
\item
  设 X 与 Y 为拓扑空间，A 为 X 的非空子空间，B 为 Y 的子空间。设 \(f: X \to Y\) 为连续映射，满足 \(f(A) \subset B\)。用 \(\varphi: X/A \to Y/B\) 表示由 f 在商上诱导的连续映射。若 f 定义了偶 (X, A) 到偶 (Y, B) 上的同伦等价，则映射 \(\varphi\) 是偶 (X/A, \(s_{X/A}\)) 到偶 (Y/B, \(s_{Y/B}\)) 上的同伦等价。设 P 为缩化为单独一点的拓扑空间，\(\alpha\) 与 \(\beta\) 为从 A 与 B 到 P 中的常值映射。注意到映射 \(\varphi\) 可等同于由 f 与 P 的恒同映射诱导的、从 \(P \cup_{\alpha} X\) 到 \(P \cup_{\beta} Y\) 中的映射。断言便由 III, p. 249 的推论 1 得出。
\end{enumerate}

\subsection{映射的锥}

设 X 与 Y 为拓扑空间，f 为从 X 到 Y 中的连续映射。设 \(\operatorname{Cyl}(f)\) 为映射 f 的柱。用 \(\alpha_{f}\) 表示从 \(X \times I\) 到 \(\operatorname{Cyl}(f)\) 中的典范映射，用 \(f_{0}\) 表示映射 \(x \mapsto \alpha_{f}(x,0)\)（从 X 到 \(\operatorname{Cyl}(f)\) 中）；这些映射分别诱导出 \(X \times I\) 与 X 到它们在 \(\operatorname{Cyl}(f)\) 中的像上的同胚。

定义 11. --- 映射 f 的锥记作 \(\operatorname{Côn}(f)\)，它是由 \(\operatorname{Cyl}(f)\) 缩塌 \(f_{0}(\mathrm{X})\) 所得的拓扑空间。

用 \(\beta_{f}^{\prime}\colon\mathrm{Y}\to\mathrm{C}\hat{\mathrm{o}}\mathrm{n}(f)\) 表示典范映射 \(\beta_{f}\colon\mathrm{Y}\to\mathrm{C}\mathrm{y}\mathrm{l}(f)\) 与从 \(\mathrm{C}\mathrm{y}\mathrm{l}(f)\) 到 \(\mathrm{C}\hat{\mathrm{o}}\mathrm{n}(f)\) 上的典范满射的复合。映射 \(\beta_{f}^{\prime}\) 是连续的，并定义了从 Y 到 \(\mathrm{C}\hat{\mathrm{o}}\mathrm{n}(f)\) 的一个闭子集上的同胚，该闭子集称为锥底，我们据此把它与 Y 等同。

用 \(\alpha_{f}^{\prime}\colon\mathrm{X}\times\mathbf{I}\to\mathrm{C}\hat{\mathrm{o}}\mathrm{n}(f)\) 表示映射 \(\alpha_{f}\) 与从 \(\mathrm{Cyl}(f)\) 到 \(\mathrm{C}\hat{\mathrm{o}}\mathrm{n}(f)\) 上的典范满射的复合。映射 \(\alpha_{f}^{\prime}\) 是一个同伦，其初始映射是常值映射，末项映射是映射 \(\beta_{f}^{\prime}\circ f\)。

若 X 为空，则映射 \(\beta_{f}^{\prime}\) 是从 Y 到 \(\operatorname{Côn}(f)\) 上的同胚。

假设 X 非空。用 s 表示空间 \(\mathrm{Cyl}(f)/f_{0}(\mathrm{X})\) 的基点；它称为锥 \(\mathrm{Côn}(f)\) 的锥顶点。由于 \(f_{0}(\mathrm{X})\) 在 \(\mathrm{Cyl}(f)\) 中闭，典范映射 \(\pi\colon\mathrm{Cyl}(f)\to\mathrm{Côn}(f)\) 诱导出 \(\mathrm{Cyl}(f)\setminus f_{0}(\mathrm{X})\) 到 \(\mathrm{Côn}(f)\setminus\{s\}\) 上的同胚（III, p. 252）。锥 \(\mathrm{Côn}(f)\) 的锥顶点不属于其锥底；Y 到 \(\mathrm{Côn}(f)\setminus\{s\}\) 中的典范单射是同伦等价（III, p. 239, 命题 6）。

设 \(\sigma_f\colon \mathrm{Cyl}(f)\times \mathbf{I}\to \mathrm{Cyl}(f)\) 为 \(f\) 的柱到其柱底上的典范收缩。对 \(c\in \mathrm{Cyl}(f)\setminus f_0(\mathrm{X})\) 与 \(t\in \mathbf{I}\)，有 \(\sigma_f(c,t)\notin f_0(\mathrm{X})\)。因此，存在唯一的映射 \(\sigma_f'\colon (\mathrm{Côn}(f)\setminus\{s\})\times \mathbf{I}\to \mathrm{Côn}(f)\setminus\{s\}\)，满足 \(\sigma_f'(\pi(c),t)=\pi(\sigma_f(c,t))\)（对 \(c\in \mathrm{Cyl}(f)\setminus f_0(\mathrm{X})\) 与 \(t\in \mathbf{I}\)）。它是连续的，并且是 \(\mathrm{Côn}(f)\setminus\{s\}\) 到 Y 上的一个强收缩。它称为 \(\mathrm{Côn}(f)\setminus\{s\}\) 到其锥底上的典范收缩。它的末项映射是 \(\mathrm{Côn}(f)\setminus\{s\}\) 到 Y 上的一个收缩映射，称为去掉锥顶点的锥到其锥底上的典范收缩映射。

命题 11（锥的泛性质）

设 Z 为拓扑空间，\(\beta: \mathrm{Y} \to \mathrm{Z}\) 为连续映射，\(\alpha: \mathrm{X} \times \mathbf{I} \to \mathrm{Z}\) 为一同伦，其初始映射是常值映射，末项映射等于 \(\beta \circ f\)。存在唯一的连续映射 \(\varphi\)，它从 \(\operatorname{Côn}(f)\) 到 Z 中，满足 \(\alpha = \varphi \circ \alpha_f'\) 且 \(\beta = \varphi \circ \beta_f'\)。

由柱的泛性质（III, p. 238, 命题 4），存在唯一的连续映射 \(h \colon \operatorname{Cyl}(f) \to \mathbf{Z}\)，满足 \(\alpha = h \circ \alpha_f\) 且 \(\beta = h \circ \beta_f\)。由于 \(\alpha\) 的初始映射是常值映射，\(h\) 在子空间 \(\alpha_f(\mathbf{X} \times \{0\})\) 上的限制是常值的。因此存在唯一的连续映射 \(\varphi \colon \operatorname{Côn}(f) \to \mathbf{Z}\)，满足 \(h = \varphi \circ \pi\)，其中 \(\pi\) 表示从 \(\operatorname{Cyl}(f)\) 到 \(\operatorname{Côn}(f)\) 上的典范满射。映射 \(\varphi\) 满足 \(\alpha = \varphi \circ \alpha_f'\) 与 \(\beta = \varphi \circ \beta_f'\)，并且是唯一具有这些性质的映射，因为 \(\alpha_f'\) 与 \(\beta_f'\) 的像覆盖 \(\operatorname{Côn}(f)\)。

例 1. --- 设 X 为拓扑空间。我们把映射 IdX 的锥称为空间 X 的锥，记作 C(X)；它是由 X × I 缩塌 X × \{0\} 所得的空间。设 \(\pi\) 为从 X × I 到 C(X) 中的典范映射；X × {[}0, 1{[} 在 \(\pi\) 下的像 C'(X) 称为空间 X 的开锥。

假设 X 非空。那么锥 C(X) 非空，其基点仍称为锥的顶点。

映射 \(((x,t),u)\mapsto(x,t(1-u))\)（从 \((\mathrm{X}\times\mathbf{I})\times\mathbf{I}\) 到 \(X\times I\) 中）是连接 \(X\times I\) 的恒同映射与映射 \((x,t)\mapsto(x,0)\) 的同伦。由此，通过取商，我们得到映射 \(\sigma\colon\mathrm{C}(\mathrm{X})\times\mathbf{I}\to\mathrm{C}(\mathrm{X})\)，满足 \(\sigma(\pi(x,t),u)=\pi(x,t(1-u))\)（对 \(x\in X\)，\(t,u\in I\)）。由 I, p. 19, 命题 10，映射 \(\sigma\) 是连续的。于是映射 \(\sigma\) 是 s 处的带基点同伦，连接 \(\mathrm{C}(\mathrm{X})\) 的恒同映射与以 \(\{s\}\) 为像的常值映射。因此，锥 \(\mathrm{C}(\mathrm{X})\) 在其顶点 s 处可缩。由于 \(\sigma(\mathrm{C}'(\mathrm{X})\times\mathbf{I})\) 包含于 \(\mathrm{C}'(\mathrm{X})\)，映射 \(\sigma\) 通过取子空间定义了 s 处的带基点同伦，连接 \(\mathrm{C}'(\mathrm{X})\) 的恒同映射与以 \(\{s\}\) 为像的常值映射；因此开锥 \(\mathrm{C}'(\mathrm{X})\) 在 s 处可缩。

注记。--- 1) 设 X 为拓扑空间，A 为 X 的子空间；用 i 表示 A 到 X 中的典范单射。典范收缩映射 \(r_{i}\)（它把柱 \(\mathrm{Cyl}(i)\) 收缩到其柱底上）把子空间 \(\alpha_{i}(\mathrm{A} \times \{0\})\) 映入 A 中。设 \(\rho: \mathrm{C}\hat{\mathrm{o}}\mathrm{n}(i) \to \mathrm{X}/\mathrm{A}\) 为由此通过取商诱导的连续映射。假设偶 (X, A) 具有同伦扩张性质。由 III, p. 243, 定理 1，映射 \(r_{i}\) 定义了偶 \((\mathrm{Cyl}(i), \alpha_{i}(\mathrm{A} \times \{0\}))\) 到偶 (X, A) 上的同伦等价。由 III, p. 253 的注 4，映射 \(\rho\) 便是偶 \((\mathrm{C}\hat{\mathrm{o}}\mathrm{n}(i), s)\) 到偶 \((\mathrm{X}, s_{\mathrm{X}/\mathrm{A}})\) 上的同伦等价。特别地，它是一个同伦等价。

\begin{enumerate}
\def\labelenumi{\arabic{enumi})}
\setcounter{enumi}{1}
\tightlist
\item
  设 X 与 Y 为拓扑空间，\(f\colon X \to Y\) 为连续映射。典范映射 \(\alpha_{f}^{\prime}\colon X \times I \to \operatorname{Côn}(f)\) 在 \(X \times \{0\}\) 上是常值的，因此定义了连续映射 \(\gamma_{f}\colon C(X) \to \operatorname{Côn}(f)\)。限制到 \(C'(X)\) 上时，映射 \(\gamma_{f}\) 是单射且严格的，并通过取子空间诱导出开锥 \(C'(X)\) 到 Y 在 \(\operatorname{Côn}(f)\) 中的补集上的同胚。
\end{enumerate}

把锥 \(C(X)\) 的锥底与空间 X 等同，并用 \(u\) 表示从 \(Y \cup_f C(X)\) 到 \(\operatorname{Côn}(f)\) 中的、由映射 \(\beta_f':Y\to\operatorname{Côn}(f)\) 与 \(\gamma_f:C(X)\to\operatorname{Côn}(f)\) 诱导的连续映射。反之，设 \(v:\operatorname{Côn}(f)\to Y\cup_f C(X)\) 为唯一的连续映射，满足 \(v\circ\alpha_f'\) 是从 \(X\times I\) 到 \(C(X)\) 上的典范映射，\(v\circ\beta_f'\) 是从 \(Y\) 到 \(Y\cup_f C(X)\) 的典范映射。映射 \(u\) 与 \(v\) 是互逆的同胚。

\section{同伦与道路}

\subsection{道路}

定义 1. --- 设 X 为拓扑空间。X 中的道路是指从 I 到 X 中的任一连续映射 c。点 \(c(0)\) 称为道路 c 的起点，点 \(c(1)\) 称为道路 c 的终点。X 中的圈是指起点与终点相等的 X 中的道路。

设 \(x\) 与 \(y\) 为 X 的点。称 X 中的道路 \(c\) 连接 \(x\) 与 \(y\)，如果 \(x\) 是它的起点且 \(y\) 是它的终点。

定义 2. --- 设 X 为拓扑空间。若 \(c(1) = d(0)\)，则称 X 中的道路 c 与 d 是可拼接的。这时把由下列公式定义的道路 \(c*d\) 称为 c 与 d 的拼接道路：

\[
(c * d) (t) = \left\{ \begin{array}{l l} c (2 t) & \text { for } 0 \leqslant t \leqslant 1 / 2, \\ d (2 t - 1) & \text { for } 1 / 2 \leqslant t \leqslant 1. \end{array} \right.\tag{1}
\]

它的起点是 c 的起点，它的终点是 d 的终点。

设 c 为 X 中的道路；把 c 的逆道路记作 \(\overline{c}\)，它是由 \(\overline{c}(t) = c(1 - t)\)（对 \(t \in I\)）定义的道路。

若 c 与 d 是 X 中两条可拼接的道路，则 \(\overline{d}\) 与 \(\overline{c}\) 也可拼接，且有 \(\overline{c*d} = \overline{d} * \overline{c}\)。对 X 中每条道路 c，有 \(\overline{\overline{c}} = c\)。

注记。--- 1) 设 X 为拓扑空间，P 为缩化为一点的拓扑空间。借助投影 pr₂ 把 P × I 与 I 等同。从 P 到 X 的映射（必然连续）之间的同伦组成的集合 C(P × I; X) 这时等同于 X 中的道路组成的集合 C(I; X)。连接以 x 为像的常值映射与以 y 为像的常值映射的同伦，对应于以 x 为起点、y 为终点的道路。上述等同与拼接及取逆道路的概念相容（参见 III, p. 230）。

\begin{enumerate}
\def\labelenumi{\arabic{enumi})}
\setcounter{enumi}{1}
\tightlist
\item
  设 X 与 Y 为拓扑空间。典范映射
\end{enumerate}

\[
\mathcal {C} (\mathrm{X} \times \mathbf {I}; \mathrm{Y}) \rightarrow \mathcal {C} (\mathbf {I}; \mathcal {C} _ {c} (\mathrm{X}; \mathrm{Y}))
\]

把每个同伦 \(\sigma\colon X \times I \to Y\)（它连接从 X 到 Y 的两个映射 \(f\) 与 \(g\)）对应到一条以 \(f\) 为起点、\(g\) 为终点的、空间 \(\mathcal{C}_c(X;Y)\) 中的道路。若 \(X\) 是局部紧空间，这个典范映射是双射（TG, X, p. 28, 定理 3）。

设 X 为拓扑空间。把 X 的道路空间记作 \(\Lambda(\mathrm{X})\)，它就是拓扑空间 \(\mathcal{C}_{c}(\mathbf{I};\mathrm{X})\)，其元素是 X 中的道路，其拓扑是紧收敛拓扑（参见 TG, X, p. 27）。若 x 是 X 的一点，用 \(\Lambda_{x}(\mathrm{X})\) 表示 \(\Lambda(\mathrm{X})\) 中由以 x 为起点的道路组成的子空间。若 y 是 X 的另一点，用 \(\Lambda_{x,y}(\mathrm{X})\) 表示 \(\Lambda(\mathrm{X})\) 中由以 x 为起点、y 为终点的道路组成的子空间。

命题 1. --- 设 X 与 Z 为拓扑空间。为使映射 \(\varphi\)（从 Z 到道路空间 \(\mathcal{C}_c(\mathbf{I};\mathrm{X})\) 中）连续，必须而且只需由 \((z,t)\mapsto \varphi (z)(t)\) 给出的映射是从 \(\mathbf{Z}\times \mathbf{I}\) 到 X 中的连续映射。特别地，映射 \((c,t)\mapsto c(t)\)（从 \(\mathcal{C}_c(\mathbf{I};\mathrm{X})\times \mathbf{I}\) 到 X 中）是连续的。

由于拓扑空间 I 是局部紧的，本命题由 TG, X, p. 28, 定理 3 得出。

把映射 \(c \mapsto c(0)\)（从 \(\Lambda(\mathrm{X})\) 到 X 中）记作 o，称为起点映射；把映射 \(c \mapsto c(1)\)（从 \(\Lambda(\mathrm{X})\) 到 X 中）记作 e，称为终点映射。映射 o 与 e 是连续的（TG, X, p. 27, 注 1）。X 中可拼接道路的对就是 X-空间 \((\Lambda(\mathrm{X}), e)\) 与 \((\Lambda(\mathrm{X}), o)\) 的纤维积的元素。

命题 2. --- 设 X 为拓扑空间。把道路 c 对应到逆道路 \(\overline{c}\) 的映射是 \(\Lambda(X)\) 到其自身上的同胚。道路的拼接 \((c,d)\mapsto c*d\) 是空间 \((\Lambda(X),e)\times_{X}(\Lambda(X),o)\) 到 \(\Lambda(X)\) 上的同胚。

映射 \(t \mapsto 1 - t\) 是空间 I 到其自身上的同胚；第一个断言由此得出。

用 C 表示纤维积 \((\Lambda(\mathrm{X}),e)\times_{\mathrm{X}}(\Lambda(\mathrm{X}),o)\)。用 \(\gamma:C\to\Lambda(X)\) 表示映射 \((c,d)\mapsto c*d\)，用 \(\delta:C\times I\to X\) 表示映射 \(((c,d),t)\mapsto(c*d)(t)\)。由公式 (1) 及命题 1，映射 \(\delta\) 在 \(C\times[0,\frac{1}{2}]\) 与 \(C\times[\frac{1}{2},1]\) 上的限制是连续的。因此映射 \(\delta\) 是连续的（TG, I, p. 19, 命题 4），映射 \(\gamma\) 亦然（命题 1）。我们有 \(c(t)=(c*d)(\frac{t}{2})\) 与 \(d(t)=(c*d)(\frac{1+t}{2})\)，因此 \(\gamma\) 是单射。

设 \(g\) 为 \(\mathbf{X}\) 中的道路；对 \(t \in \mathbf{I}\)，令

\[
c _ {g} (t) = g \bigl (\frac {t}{2} \bigr) \quad \text{且} \quad d _ {g} (t) = g \bigl (\frac {1 + t}{2} \bigr).
\]

这样定义的映射 \(c_{g}\) 与 \(d_{g}\) 是 X 中可拼接的道路，且有 \(c_{g} * d_{g} = g\)。此外，映射 \(g \mapsto c_{g}\) 与 \(g \mapsto d_{g}\) 是空间 \(\Lambda(\mathrm{X})\) 到其自身的连续映射（命题 1）。由此推出映射 \(\gamma\) 是同胚。

推论. --- 设 X 为拓扑空间，x、y、z 为 X 的点。映射 \(c \mapsto \overline{c}\)（从 \(\Lambda_{x,y}(X)\) 到 \(\Lambda_{y,x}(X)\)）与映射 \((c,d) \mapsto c * d\)（从 \(\Lambda_{x,y}(X) \times \Lambda_{y,z}(X)\) 到 \(\Lambda_{x,z}(X)\)）都是连续的。

\subsection{道路连通空间}

定义 3. --- 若对拓扑空间 X 的每一对点 \((x,y)\)，都存在以 x 为起点、y 为终点的道路，则称 X 是道路连通的。

例 1. --- R 的每个区间、数值空间的每个凸子集，或更一般地，R 上拓扑向量空间的每个凸子集，都是道路连通的。

命题 3. --- 道路连通空间在连续映射下的像是道路连通的。

设 X 为道路连通空间，\(f: X \to Y\) 为连续映射。设 x 与 y 为 \(f(X)\) 的两点；设 \(x'\) 与 \(y'\) 为 X 中满足 \(f(x') = x\) 与 \(f(y') = y\) 的点。存在 X 中以 \(x'\) 为起点、\(y'\) 为终点的道路 c。于是 \(f \circ c\) 是 \(f(X)\) 中以 x 为起点、y 为终点的道路。

命题 4. --- 道路连通拓扑空间是连通的。

由于 R 的每个区间都连通，道路的像是连通的（TG, I, p. 82, 命题 4）。空空间是连通的。由于非空道路连通空间是从其中一点出发的道路的像的并，它是连通的（TG, I, p. 81, 命题 2）。

注意到（III, p. 256, 注 1）X 中的道路与从缩化为一点的拓扑空间 P 到 X 中的映射之间的同伦的等同，III, p. 230 的命题 1 蕴涵：

命题 5. --- 在拓扑空间中，关系“存在以 x 为起点、y 为终点的道路”是一个等价关系。

拓扑空间 X 的道路连通分支指上述关系的等价类。道路连通分支是道路连通空间。设 x 为 X 的一点。x 的道路连通分支是 X 中包含 x 的道路连通子空间的并。它也是 X 中以 x 为起点的道路的像的并。

例。--- 2) 交非空的一族道路连通集的并是道路连通的（参见 TG, I, p. 81, 命题 2）。

\begin{enumerate}
\def\labelenumi{\arabic{enumi})}
\setcounter{enumi}{2}
\tightlist
\item
  数值空间（或更一般地，R 上的拓扑向量空间）的在其中一点处为星形的子集是道路连通的。
\end{enumerate}

设 X 为拓扑空间。用 \(\pi_{0}(X)\) 表示 X 的道路连通分支组成的集合。设 P 为缩化为一点的拓扑空间。从 X 到 {[}P; X{]} 的映射把 X 的每一点 x 对应到同伦类 \(\varphi_{x}\)，即像为 x 的映射 \(f_{x}: P \to X\) 的同伦类，它通过取商定义了一个从 \(\pi_{0}(X)\) 到 {[}P; X{]} 上的双射，称为典范双射。我们有 \(\pi_{0}(\varnothing) = \varnothing\)。为使非空拓扑空间道路连通，必须而且只需 \(\pi_{0}(X)\) 是单元素集。

设 X 与 Y 为拓扑空间，\(f\colon\mathrm{X}\to\mathrm{Y}\) 为连续映射。X 的每个道路连通分支 C 的像 \(f(\mathrm{C})\) 是道路连通的（III, p. 258, 命题 3），因此包含于 Y 的某个道路连通分支。用 \(\pi_{0}(f)\colon\pi_{0}(\mathrm{X})\to\pi_{0}(\mathrm{Y})\) 表示这样的映射：它把 X 的道路连通分支 C 对应到满足 \(f(\mathrm{C})\subset\mathrm{C}'\) 的唯一的 Y 的道路连通分支 C'。若把 \(\pi_0(\mathrm{X})\) 与 \(\pi_0(\mathrm{Y})\) 分别等同于 {[}P; X{]} 与 {[}P; Y{]}，则 \(\pi_0(f)\) 等同于映射 \(\chi \mapsto [f] \circ \chi\)（从 {[}P; X{]} 到 {[}P; Y{]}）。特别地，若 \(f\) 与 \(g\) 是从 X 到 Y 的互相同伦的映射，则有 \(\pi_0(f) = \pi_0(g)\)（III, p. 230, 命题 2）。

设 Z 为拓扑空间，\(g: Y \to Z\) 为连续映射；我们有 \(\pi_{0}(g \circ f) = \pi_{0}(g) \circ \pi_{0}(f)\)。特别地，若 Z = X 且 f 与 g 互为同伦逆，则有 \(\pi_{0}(g) \circ \pi_{0}(f) = \pi_{0}(\mathrm{Id}_{\mathrm{X}}) = \mathrm{Id}_{\pi_{0}(\mathrm{X})}\) 与 \(\pi_{0}(f) \circ \pi_{0}(g) = \pi_{0}(\mathrm{Id}_{\mathrm{Y}}) = \mathrm{Id}_{\pi_{0}(\mathrm{Y})}\)，这证明 \(\pi_{0}(f)\) 与 \(\pi_{0}(g)\) 是互逆的双射。因此，与道路连通空间同伦等价的空间本身是道路连通的。特别地，与一点同伦等价的空间是道路连通的。

命题 6. --- 设 \((\mathrm{Y}_{j})_{j\in\mathrm{J}}\) 为一族拓扑空间。映射

\[
(\pi_ {0} (\mathrm{pr} _ {j})) \colon \pi_ {0} (\prod_ {j \in \mathrm{J}} \mathrm{Y} _ {j}) \to \prod_ {j \in \mathrm{J}} \pi_ {0} (\mathrm{Y} _ {j})
\]

是双射。特别地，一族道路连通空间的积空间是道路连通的。

这由 III, p. 231 的命题 3 得出，其中取 X 为缩化为单独一点的空间。

\subsection{局部道路连通空间}

定义 4. --- 若拓扑空间的每一点都有由道路连通邻域组成的基本邻域系，则称它是局部道路连通的。

命题 7. --- 局部道路连通拓扑空间是局部连通的。其连通分支与道路连通分支重合。特别地，若局部道路连通拓扑空间连通，则它道路连通。

第一个断言由命题 4 得出。我们来证明第二个断言。设 X 为局部道路连通拓扑空间，C 为 X 的一个道路连通分支。C 的每一点都有道路连通的邻域，它包含于 C；因此 C 是 X 的开子集。由于诸道路连通分支构成 X 的一个分割，这样的分支也是闭的。由于它连通（III, p. 258, 命题 4），它是一个连通分支（TG, I, p. 83）。第三个断言由第二个得出。

于是注意，说一个拓扑空间连通且局部道路连通并无歧义。

推论 1. --- 局部道路连通拓扑空间的每个连通开子集是道路连通的。

推论 2. --- 设 B 为局部道路连通拓扑空间，E 为平展 B-空间。空间 E 是局部道路连通的。若它连通，则它道路连通。

第一个断言由定义 4 及平展态射的定义（I, p. 28, 定义 2）立即得出。第二个由命题 7 从它推出。

为使拓扑空间 X 局部道路连通，必须而且只需 X 的开子集的每个道路连通分支都是 X 的开子集（参见 I, p. 85, 命题 11 的证明）。若空间 X 局部道路连通，则 X 的每一点都有由开的道路连通邻域组成的基本邻域系。

命题 8. --- 局部道路连通空间的每个商空间都是局部道路连通的。

设 X 为局部道路连通空间，R 为 X 上的等价关系，\(\varphi\colon\mathrm{X}\to\mathrm{X}/\mathrm{R}\) 为典范映射。只需证明 X/R 的开子集的道路连通分支都是开的。设 A 为 X/R 的开子集，C 为 A 的一个道路连通分支。设 \(x\in\bar{\varphi}^{1}(\mathrm{C})\)，K 为 x 在 \(\bar{\varphi}^{1}(\mathrm{A})\) 中的道路连通分支。集合 \(\varphi(\mathrm{K})\) 是道路连通的（III, p. 258, 命题 3），包含于 A 且包含 \(\varphi(x)\)；于是有 \(\varphi(\mathrm{K})\subset\mathrm{C}\)，从而 \(\mathrm{K}\subset\bar{\varphi}^{1}(\mathrm{C})\)。这证明 \(\bar{\varphi}^{1}(\mathrm{C})\) 是 \(\bar{\varphi}^{1}(\mathrm{A})\) 的道路连通分支的并。由于 X 局部道路连通且 \(\bar{\varphi}^{1}(\mathrm{A})\) 在 X 中开，\(\bar{\varphi}^{1}(\mathrm{C})\) 在 X 中开。因此 C 在 X/R 中开。命题得证。

命题 9. --- 一族局部道路连通空间（除有限个外都连通）的积是局部道路连通的。

设 \((\mathrm{X}_{j})_{j\in\mathrm{J}}\) 为一族局部道路连通空间，除有限个外都连通。设 \(x=(x_{j})_{j\in\mathrm{J}}\) 为积空间 \(X=\prod_{j\in J}X_{j}\) 的一点。按假设，x 在 X 中的基本邻域系由形如 \(V=\prod_{j\in J}V_{j}\) 的集合组成，其中对每个 \(j\in J\)，\(V_{j}\) 是 \(x_{j}\) 在 \(X_{j}\) 中的道路连通邻域，且 \(V_{j}=X_{j}\)（除有限多个指标 \(j\in J\) 外）（TG, I, p. 24）。这些集合都是道路连通的（III, p. 260, 命题 6），故 X 局部道路连通。

推论. --- 数值空间的每个开子集都是局部道路连通的。

R 的每一点都有由区间组成的邻域基；因此 R 是局部道路连通的。由命题 9，数值空间 \(R^{n}\)（对每个整数 \(n \geqslant 1\)）都是局部道路连通的。此外，局部道路连通拓扑空间的开子集仍是局部道路连通的，推论得证。

例. --- 设 G 为 \(\mathbf{GL}(n,\mathbf{R})\) 中由行列式严格正的 n 阶方阵组成的子群。群 G 是连通的且局部道路连通的。

由 A, III, p. 104, 命题 17，群 \(\mathbf{SL}(n,\mathbf{R})\) 由元素 \(B_{ij}(\lambda)\)（\(1\leqslant i,j\leqslant n\)，\(i\neq j\)，\(\lambda\in\mathbf{R}\)）生成。映射 \(\lambda\mapsto B_{ij}(\lambda)\)（从 \(\mathbf{R}\) 到 \(\mathbf{GL}(n,\mathbf{R})\) 中）是连续的，它们的像是 \(\mathbf{SL}(n,\mathbf{R})\) 的连通子集；由于它们都包含单位矩阵 \(I_n\)，它们的并是连通的（TG, I, p. 81, 命题 2）。因此群 \(\mathbf{SL}(n,\mathbf{R})\) 是连通的（TG, III, p. 8, 命题 7）。设 A 为 G 中由形如 diag(1, \(\ldots\),1, \(\lambda\))（\(\lambda\in\mathbf{R}_{+}^{*}\)）的矩阵组成的子群；它是连通的，且有 \(\mathbf{GL}(n,\mathbf{R})=\mathbf{A}\cdot\mathbf{SL}(n,\mathbf{R})\)。由此推出群 G 是连通的。由于它是 \(\mathbf{R}_{+}^{*}\) 在从 \(\mathbf{M}_{n}(\mathbf{R})\) 到 \(\mathbf{R}\) 中的行列式映射下的原像，它是 \(\mathbf{M}_{n}(\mathbf{R})\) 的开子集；因此它是局部道路连通的（前述推论），也是道路连通的（III, p. 261, 推论 1）。

命题 10. --- 设 X 为拓扑空间。映射 \((o, e)\)（从 \(\Lambda(\mathrm{X}) = \mathcal{C}_{c}(\mathbf{I}; \mathrm{X})\) 到 \(X \times X\) 中，它把 X 中的道路 c 对应到对 \((c(0), c(1))\)）是连续的。若空间 X 局部道路连通，则该映射是开的。

我们已经知道 \(\Lambda(\mathrm{X})\) 到 X 中的起点映射 o 与终点映射 e 是连续的（III, p. 257），第一个断言由此得出。

引理. --- 设 \(c\colon \mathbf{I}\to \mathrm{X}\) 为道路，W 为 \(c\) 在空间 \(\mathcal{C}_c(\mathbf{I};\mathrm{X})\) 中的邻域。存在实数 \(\varepsilon \in ]0,1 / 2]\)、邻域 \(\mathrm{V}_0\)（含 \(c(0)\)）与邻域 \(\mathrm{V}_1\)（在 X 中、含 \(c(1)\)），具有下列性质：我们有 \(c([0,\varepsilon])\subset \mathrm{V}_0\)，\(c([1 - \varepsilon,1])\subset \mathrm{V}_1\)，且每条道路 \(c^{\prime}\colon \mathbf{I}\rightarrow \mathrm{X}\)，若它满足

\[
\left\{ \begin{array}{l l} c ^ {\prime} (t) \in \mathrm{V} _ {0} & \text { for } 0 \leqslant t \leqslant \varepsilon , \\ c ^ {\prime} (t) = c (t) & \text { for } \varepsilon \leqslant t \leqslant 1 - \varepsilon , \\ c ^ {\prime} (t) \in \mathrm{V} _ {1} & \text { for } 1 - \varepsilon \leqslant t \leqslant 1, \end{array} \right.\tag{2}
\]

，便属于 W。

由紧收敛拓扑的定义（TG, X, p. 26, 定义 1），存在有限集 J、X 的开集族 \((\mathrm{U}_{j})_{j\in\mathrm{J}}\) 与 I 的紧子集族 \((\mathrm{K}_{j})_{j\in\mathrm{J}}\)，使得集合 \(W'\)（其中的道路 \(c'\) 对每个指标 j 满足 \(c'(\mathrm{K}_{j})\subset\mathrm{U}_{j}\)）是包含于 W 中的 c 的邻域。然后用 \(A_{0}\)（相应地，\(A_{1}\)）表示满足 \(0\in K_{j}\)（相应地，\(1\in K_{j}\)）的指标 j 的集合；令 \(V_{0}=\bigcap_{j\in A_{0}}U_{j}\)，\(V_{1}=\bigcap_{j\in A_{1}}U_{j}\)。

由于映射 \(c\) 连续，存在实数 \(\varepsilon \in [0,1/2]\)，使得 \(c([0,\varepsilon]) \subset \mathrm{V}_0\)，\(c([1-\varepsilon,1]) \subset \mathrm{V}_1\)，且 \([0,\varepsilon] \cap \mathrm{K}_j = \emptyset\)（对所有 \(j \notin \mathrm{A}_0\)），\([1-\varepsilon,1] \cap \mathrm{K}_j = \emptyset\)（对所有 \(j \notin \mathrm{A}_1\)）。设 \(c'\) 为满足条件 (2) 的道路。我们来证明 \(c' \in \mathrm{W}'\)。设 \(j \in \mathrm{J}\)，\(t \in \mathrm{K}_j\)。若 \(\varepsilon \leqslant t \leqslant 1-\varepsilon\)，则 \(c'(t) = c(t)\) 属于 \(\mathrm{U}_j\)。若 \(0 \leqslant t \leqslant \varepsilon\)，则 \(c'(t) \in \mathrm{V}_0\)；由 \(\varepsilon\) 的取法有 \(j \in \mathrm{A}_0\)，因此 \(c'(t) \in \mathrm{U}_j\)。类似地，若 \(1-\varepsilon \leqslant t \leqslant 1\)，则有 \(j \in \mathrm{A}_1\) 且 \(c'(t) \in \mathrm{V}_1 \subset \mathrm{U}_j\)。于是 \(c'(\mathrm{K}_j) \subset \mathrm{U}_j\)，从而 \(c'\) 属于 \(\mathrm{W}'\)，故属于 W。

现在证明命题 10 的第二个断言。设空间 X 局部道路连通。设 \(c\) 为 X 中的道路，W 为 \(c\) 在 \(\mathcal{C}_c(\mathbf{I};\mathbf{X})\) 中的邻域。取引理中的 \(\varepsilon\)、\(\mathrm{V}_0\) 与 \(\mathrm{V}_1\)。设 \(\mathrm{T}_0\) 与 \(\mathrm{T}_1\) 为 \(c(0)\) 与 \(c(1)\) 的道路连通邻域，分别包含于 \(\mathrm{V}_0\) 与 \(\mathrm{V}_1\) 中。存在实数 \(\theta\)，满足 \(0 < \theta < \varepsilon\)，且 \(c([0,\theta]) \subset \mathrm{T}_0\)，\(c([1 - \theta,1]) \subset \mathrm{T}_1\)。设 \(x_0 \in \mathrm{T}_0\)，\(x_1 \in \mathrm{T}_1\)；设 \(c_0\) 为以 \(x_0\) 为起点、\(c(\theta)\) 为终点的道路（在 \(\mathrm{T}_0\) 中），\(c_1\) 为以 \(x_1\) 为起点、\(c(1 - \theta)\) 为终点的道路

（在 \(\mathrm{T}_1\) 中）。令

\[
c ^ {\prime} (t) = \left\{ \begin{array}{l l} c _ {0} (t / \theta) & \text { for } 0 \leqslant t \leqslant \theta , \\ c (t) & \text { for } \theta \leqslant t \leqslant 1 - \theta , \\ c _ {1} ((1 - t) / \theta) & \text { for } 1 - \theta \leqslant t \leqslant 1. \end{array} \right.
\]

这样便定义了一条道路 \(c'\)，它连接 \(x_{0}\) 与 \(x_{1}\) 并满足条件 (2)。这证明 W 在映射 \((o, e)\) 下在 X × X 中的像包含邻域 \(T_{0} \times T_{1}\)（\((c(0), c(1))\) 的），命题得证。

推论. --- 设 X 为局部道路连通拓扑空间，x 为 X 的一点。映射 \(c \mapsto c(1)\)（从 \(\Lambda_{x}(\mathrm{X})\) 到 X 中）是开的。

由本命题，映射 \(\varphi\colon\Lambda(\mathrm{X})\to\mathrm{X}\times\mathrm{X}\)（由 \(\varphi(c)=(c(0),c(1))\) 定义）是开的，且有 \(\Lambda_{x}(\mathrm{X})=\overline{\varphi}^{1}(\{x\}\times\mathrm{X})\)。因此，映射 \(\Lambda_{x}(\mathrm{X})\to\{x\}\times\mathrm{X}\)（由 \(\varphi\) 诱导）是开的（TG, I, p. 30, 命题 2），它与第二投影 \(\mathrm{pr}_{2}\) 的复合亦然。

\subsection{连通性与道路连通性之间的联系}

道路连通空间是连通的（III, p. 258, 命题 4）。存在连通的、甚至局部连通的空间，它们不是道路连通的（参见 III, p. 331, 习题 2 与 4）。然而：

命题 11. --- 连通且局部连通、其拓扑可由一个使之为完备的距离定义的拓扑空间，是道路连通的。

设 X 为拓扑空间。把 X 中非空有限序列 \(\mathrm{T} = (\mathrm{W}_{i})_{1 \leqslant i \leqslant n}\)（由 X 的连通开子集组成，满足 \(W_{i} \cap W_{i+1} \neq \varnothing\)（对 \(1 \leqslant i \leqslant n-1\)））称为 X 中的列车。称 n 为列车 T 的长，称诸 \(W_{i}\) 为它的车厢。若 X 赋有与它的拓扑相容的距离，则称列车 T 的诸车厢直径的最大值为列车 T 的宽度。若点 a 属于第一节车厢而点 b 属于最后一节车厢，则称该列车连接点 a 与点 b。把对 \((\mathrm{T}', f)\)（由列车 \(\mathrm{T}' = (\mathrm{W}_{j}')_{1 \leqslant j \leqslant m}\) 与严格递增映射 \(f: \{0, 1, \ldots, n\} \to \{0, 1, \ldots, m\}\) 构成，满足 \(f(0) = 0\)，\(f(n) = m\)，且 \(W_{j}' \subset W_{i}\)（对 \(1 \leqslant i \leqslant n\) 与 \(f(i-1) < j \leqslant f(i)\)））称为 T 的加细。

引理 1. --- 设 X 为连通且局部连通的度量空间，a 与 b 为 X 的点，\(\varepsilon\) 为实数 \(>0\)。X 中存在宽度 \(\leqslant \varepsilon\) 的列车连接 a 与 b。

更精确地说，每条连接 a 与 b 的列车 T 都有加细 \((\mathrm{T}^{\prime}, f)\)，其中 \(T^{\prime}\) 是宽度 \(\leqslant\varepsilon\) 且连接 a 与 b 的列车。

关系“存在宽度 \(\leqslant\varepsilon\) 的列车连接 x 与 y”是 X 中 x 与 y 之间的等价关系。一点 x 的等价类包含 x 的每个直径 \(\leqslant\varepsilon\) 的连通开邻域，而由于 X 局部连通，x 具有这样的邻域。因此这个关系下的等价类都是开的，从而也是闭的。由于 X 连通，它们至多有一个。这就证明了第一个断言。

设 \(\mathrm{T}=(\mathrm{W}_{i})_{1\leqslant i\leqslant n}\) 为 X 中连接 a 与 b 的列车。令 \(x_{0}=a\)，\(x_{n}=b\)，并对 \(1\leqslant i\leqslant n-1\) 选取一点 \(x_{i}\)，它属于非空集 \(W_{i}\cap W_{i+1}\)。对 \(1\leqslant i\leqslant n\)，开集 \(W_{i}\) 连通且局部连通，而 \(x_{i-1}\)、\(x_{i}\) 是它的两个点；由前段，存在列车 \((\mathrm{W}_{i,k})_{1\leqslant k\leqslant m_{i}}\)（在 \(W_{i}\) 中，宽度 \(\leqslant\varepsilon\)，连接 \(x_{i-1}\) 与 \(x_{i}\)）。

令 \(m = m_{1} + \cdots + m_{n}\)。对 \(1 \leqslant j \leqslant m\)，令 \(W_{j}^{\prime} = W_{i,k}\)，其中 \((i, k)\) 是满足 \(1 \leqslant i \leqslant n\)，\(1 \leqslant k \leqslant m_{i}\) 且 \(j = m_{1} + \cdots + m_{i-1} + k\) 的唯一的整数对。对 \(0 \leqslant i \leqslant n\)，令 \(f(i) = m_{1} + \cdots + m_{i}\)。于是 \(T^{\prime} = (W_{j}^{\prime})_{1 \leqslant j \leqslant m}\) 是 X 中宽度 \(\leqslant \varepsilon\)、连接 a 与 b 的列车，且 \((T^{\prime}, f)\) 是 T 的加细。

现在证明命题 11。给连通且局部连通的空间 X 赋以与它的拓扑相容、且使之为完备的距离 d。设 a 与 b 为 X 的点。引理 1 使我们可以归纳地构造序列 \((\mathrm{T}_s)_{s \geqslant 0}\) 与 \((f_s)_{s \geqslant 1}\)，使得对每个 \(s \geqslant 0\)，\(\mathrm{T}_s = (\mathrm{W}_{s,i})_{1 \leqslant i \leqslant n_s}\) 是宽度 \(\leqslant 2^{-s}\)、连接 a 与 b 的列车，且 \((\mathrm{T}_{s+1}, f_{s+1})\) 是 \(\mathrm{T}_s\) 的加细。

我们可以归纳地选取，对 \(s \geqslant 0\)，严格递增映射 \(g_{s} \colon \{0,1,\ldots,n_{s}\} \to I\)，满足 \(g_{s}(0) = 0\) 与 \(g_{s}(n_{s}) = 1\)，使得 \(g_{s+1} \circ f_{s} = g_{s}\)。对 \(s \geqslant 0\)，用下式定义子集 \(A_{s}\)（\(I \times X\) 中的）：

\[
\mathrm{A} _ {s} = \bigcup_ {1 \leqslant i \leqslant n _ {s}} ([ g _ {s} (i - 1), g _ {s} (i) ] \times \mathrm{W} _ {s, i}).
\]

序列 \((\mathrm{A}_{s})_{s\geqslant0}\) 是递减的：事实上，对每个整数 \(j\in\{1,\ldots,n_{s+1}\}\)，存在唯一的整数 \(i\in\{1,\ldots,n_{s}\}\) 满足 \(f_{s}(i-1)<j\leqslant f_{s}(i)\)，且有

\[
[ g _ {s + 1} (j - 1), g _ {s + 1} (j) ] \subset [ g _ {s} (i - 1), g _ {s} (i) ], \quad \mathrm{W} _ {s + 1, j} \subset \mathrm{W} _ {s, i}.
\]

设 \(t \in \mathbf{I}\)。对所有 \(s \geqslant 0\)，用 \(\mathrm{A}_s(t)\) 表示由 \(x \in \mathrm{X}\)（满足 \((t, x) \in \mathrm{A}_s\)）组成的集合。集合 \(\mathrm{A}_s(t)\) 或者是列车 \(\mathrm{T}_s\) 的某一节车厢，或者是它的两节相邻车厢的并；因此它是 \(\mathrm{X}\) 的非空子集，直径 \(\leqslant 2^{1-s}\)。序列 \((\mathrm{A}_s(t))_{s \geqslant 0}\) 是递减的。它的项组成的集合是一个柯西滤子基。它收敛于一点 \(c(t)\)（因为度量空间 \(\mathrm{X}\) 是完备的）。

由于 \(a\) 属于每个集合 \(\mathrm{A}_s(0) = \mathrm{W}_{s,1}\)，我们有 \(c(0) = a\)；类似地有 \(c(1) = b\)。设 \(t \in \mathbf{I}\)，并设 \(s\) 为整数 \(\geqslant 0\)。点 \(t\) 在 \(\mathbf{I}\) 中有下列形式之一的邻域 V：\([g_s(0), g_s(1)[\)，\(]g_s(i - 1), g_s(i + 1)[\)（\(i\) 为满足 \(1 \leqslant i \leqslant n_s - 1\) 的整数），或 \(]g_s(n_s - 1), g_s(n_s)]\)。视情形而定，集合 \(c(\mathrm{V})\) 包含于列车 \(\mathrm{T}_s\) 的第一节车厢、两节相邻车厢的并、或最后一节车厢的闭包中。因此它的直径 \(\leqslant 2^{1-s}\)，且 \(d(c(t), c(t')) \leqslant 2^{1-s}\)（对 \(t' \in \mathrm{V}\)）。这证明映射 \(c\colon \mathbf{I} \to \mathbf{X}\) 是连续的。于是 \(c\) 是 X 中连接 \(a\) 与 \(b\) 的道路，故 X 道路连通。

推论 1. --- 其拓扑可由一个使之为完备的距离定义的局部连通拓扑空间，是局部道路连通的。

设 X 为这样的空间，U 为 X 的开且连通的子集。由下面的引理 2，U 上存在与它的拓扑相容、且使之为完备的距离。由于 U 局部连通，由命题 11 推出 U 是道路连通的。

引理 2. --- 设 X 为完备度量空间，U 为 X 的开子集。U 上存在与它的拓扑相容、且使之为完备的距离。

我们将采用 TG, IX, p. 57 命题 2 的证明中的论证。可以假设 U 与 X 不同。设 V 为乘积 \(\mathbf{R} \times \mathbf{X}\) 的由点 \((t,x)\)（满足 \(td(x,\mathrm{X}-\mathrm{U})=1\)）组成的子集；\(\mathbf{R} \times \mathbf{X}\) 的子空间 V 是闭的，而从 V 到 U 的映射 \((t,x) \mapsto x\) 是同胚（TG, IX, p. 13, 命题 3）。\(\mathbf{R} \times \mathbf{X}\) 上存在距离 \(d'\)（与它的拓扑相容，且使 \(\mathbf{R} \times \mathbf{X}\) 为完备）（TG, IX, p. 15, 推论 2 及 TG, II, p. 17, 命题 10）。空间 V 对由 \(d'\) 诱导的距离是完备的（TG, II, p. 16, 命题 8），引理得证。

推论 2. --- 局部紧、局部连通且可度量化的拓扑空间是局部道路连通的。若它连通，则它道路连通。

只需证明第一个断言（III, p. 260, 命题 7）。设 X 为局部紧且局部连通的度量空间。X 的开、连通且相对紧的子集构成 X 的拓扑的基。设 U 为这样一个集合；由于 U 在它的闭包（紧从而完备的度量空间）中是开的，由引理 2，U 上存在与它的拓扑相容、且使之为完备的距离。由 III, p. 264 的命题 11 推出 U 是道路连通的，推论得证。

\subsection{道路连续映射}

定义 5. --- 设 X 与 Y 为拓扑空间。映射 \(f: X \to Y\) 称为道路连续的，如果对 X 中的每条道路 \(c\)，映射 \(f \circ c: I \to Y\) 都连续。

注. --- 设空间 X 道路连通。设 x 为 X 的一点。为使 f 道路连续，只需对每条以 x 为起点的道路 c，映射 \(f \circ c\) 连续。事实上，设 d 为 X 中任意道路，c 为 X 中以 x 为起点、\(d(0)\) 为终点的道路。若映射 \(f \circ (c * d)\) 连续，则映射 \(f \circ d\) 也连续，因为有 \(f \circ d(t) = f \circ (c * d)((t + 1)/2)\)。

连续映射是道路连续的。反之不总成立；下面的命题 12 与 13 给出了断言道路连续映射连续的判据。

命题 12. --- 设 X 与 Y 为拓扑空间，\(f: \mathrm{X} \to \mathrm{Y}\) 为映射。假设空间 X 局部道路连通，且 X 的每一点都有可数的基本邻域系。若映射 \(f\) 是道路连续的，则它连续。

由（TG, IX, p. 18），只需证明：对 X 的每一点 \(x\) 与每个序列 \((x_n)_{n \geqslant 1}\)（由 X 中趋于 \(x\) 的点组成），序列 \((f(x_n))_{n \geqslant 1}\) 趋于 \(f(x)\)。删去序列 \((x_n)_{n \geqslant 1}\) 的开头若干项，可化归为序列的各项都属于 \(x\) 在 X 中的同一道路连通邻域的情形。由下面的引理，于是存在道路 \(c: \mathbf{I} \to \mathbf{X}\)，满足 \(c(0) = x\) 与 \(c(1/n) = x_n\)（对 \(n \geqslant 1\)）。若映射 \(f\) 道路连续，则映射 \(f \circ c\) 连续，而元素 \(f(x) = f(c(0))\) 是序列 \((f(c(1/n)))_{n \geqslant 1}\) 即序列 \((f(x_n))_{n \geqslant 1}\) 的极限，推论得证。

引理. --- 设 X 为连通且局部道路连通的拓扑空间，x 为 X 的一点。假设点 x 有可数的基本邻域系。那么，对每个趋于 x 的 X 的点的序列 \((x_{n})_{n\geqslant1}\)，存在 X 中的道路 c，满足 \(c(0)=x\) 与 \(c(1/n)=x_{n}\)（对 \(n\geqslant1\)）。

设 \((\mathrm{W}_{m})_{m\geqslant1}\) 为 x 的基本邻域系。令 \(V_{0}=X\)，且对每个 \(m\geqslant1\)，设 \(V_{m}\) 为包含于 \(V_{m-1}\cap W_{m}\) 中的 x 的道路连通邻域。

对每个整数 \(n \geqslant 1\)，用 \(m_n\) 表示最大整数 \(m \leqslant n\)，它满足 \(x_k \in V_m\)（对所有 \(k \geqslant n\)）。序列 \((m_n)_{n \geqslant 1}\) 按定义是递增的；由于序列 \((x_n)_{n \geqslant 1}\) 趋于 \(x\)，它趋于无穷。对每个整数 \(n \geqslant 1\)，设 \(c_n: \mathbf{I} \to V_{m_n}\) 为以 \(x_{n+1}\) 为起点、\(x_n\) 为终点的道路（在 \(V_{m_n}\) 中）。如下定义映射 \(c: \mathbf{I} \to X\)：令 \(c(0) = x\)，且 \(c(t) = c_n(n(n+1)t - n)\)（当 \(1/(n+1) < t \leqslant 1/n\) 时）。有 \(c(1/n) = x_n\)（对所有 \(n \geqslant 1\)）。因此映射 \(c\) 在每个形如 \([1/(n+1), 1/n]\)（\(n \geqslant 1\)）的区间上连续，从而在区间 \([0,1]\) 上连续。若 \(t \leqslant 1/n\)，则点 \(c(t)\) 属于 \(V_{m_n}\)；因此映射 \(c\) 在 0 处连续。

命题 13. --- 设 \(p \colon \mathrm{E} \to \mathrm{B}\) 为分离的平展态射，\(s \colon \mathrm{B} \to \mathrm{E}\) 为 \(p\) 的截面。若空间 B 局部道路连通且截面 \(s\) 道路连续，则它连续。

设 \(b\) 为 B 的一点；我们来证明 \(s\) 在点 \(b\) 处连续。由于 \(p\) 是平展映射，存在 \(b\) 的邻域 V 与定义在 V 中的连续局部截面 \(s'\)（\(p\) 的），满足 \(s'(b) = s(b)\)（I, p. 33, 命题 9）。由于 B 局部道路连通，可以假设 V 道路连通。对 V 中每条道路 \(c\)（以 \(b\) 为起点），映射 \(s \circ c\) 与 \(s' \circ c\) 是映射 \(c \colon \mathbf{I} \to \mathbf{B}\) 到 E 中的两个连续提升，且有 \(s \circ c(0) = s' \circ c(0) = s(b)\)。由 I, p. 34 的推论 1，有 \(s \circ c = s' \circ c\)，特别地有 \(s \circ c(1) = s' \circ c(1)\)。由于 V 的每一点都是 V 中从 \(b\) 出发的道路的终点，映射 \(s\) 与 \(s'\) 在 V 上一致。因此映射 \(s\) 在 V 中连续。

推论. --- 设 B 为拓扑空间，(E, p) 与 (E', p') 为两个 B-空间。假设映射 p 是平展且分离的，空间 E' 局部道路连通。那么每个道路连续映射 \(f: E' \to E\)（满足 \(p \circ f = p'\)）都是连续的。

映射 \(pr_{1}: E'\times_{B} E \rightarrow E'\) 是平展且分离的（I, p. 31, 命题 8 及 I, p. 27, 命题 4），而映射 \(x \mapsto (x, f(x))\) 是道路连续的截面。由命题 13，它连续，因此 f 连续。

\subsection{关于紧可度量化拓扑空间的补充结果}

我们给二元素集 \(\{0,1\}\) 赋以离散拓扑，给集合 \(\{0,1\}^{N}\) 赋以乘积拓扑。拓扑空间 \(\{0,1\}^{N}\) 是紧的（TG, I, p. 63, 定理 3）、可度量化的（TG, IX, p. 15, 推论 2）、非空的、全不连通的（TG, I, p. 84, 命题 10），且没有孤立点。

命题 14. --- 每个紧、可度量化、非空、全不连通且无孤立点的拓扑空间都同胚于 \(\{0,1\}^{N}\)。

设 X 为这样的拓扑空间。给它赋以与它的拓扑相容的距离。由于空间 X 紧，它对这个距离是完备的（TG, II, p. 27, 定理 1）。

引理 3. --- 设 \(\varepsilon\) 为实数 \(>0\)。存在整数 \(m \geqslant 1\)，使得对每个整数 \(n \geqslant m\)，X 都有一个由 \(n\) 个直径 \(\leqslant \varepsilon\) 的非空开闭集组成的分割。

X 的每一点都有直径 \(\leqslant\varepsilon\) 的开闭邻域（TG, II, p. 32, 命题 6 的推论）。由于空间 X 紧，它有一个由这样的集合组成的有限覆盖。选取其中之一 \((\mathrm{U}_{i})_{1\leqslant i\leqslant m}\)，使 m 最小。由于 X 非空，有 \(m\geqslant1\)。对 \(1\leqslant i\leqslant m\)，用 \(V_{i}\) 表示 \(U_{i}\) 与

\(X\setminus U_k\)（\(k < i\)）的交。于是 \((V_i)_{1 \leqslant i \leqslant m}\) 是 X 的一个分割，由 m 个直径 \(\leqslant \varepsilon\) 的非空开闭集组成。

由于 X 没有孤立点，X 的每个非空开闭子集 V 都至少含两点。此外，由于 X 紧且全不连通，V 是两个不交的非空开闭子集的并（同前所引）。由此归纳地得出，对每个整数 \(n \geqslant m\)，X 都有一个由 n 个直径 \(\leqslant \varepsilon\) 的非空开闭集组成的分割。

现在来完成命题 14 的证明。X 的每个非空开闭子空间都是紧、全不连通、无孤立点的度量空间。于是引理 3 使我们可以归纳地构造有限集的序列 \((\mathrm{J}_{n})_{n\geqslant0}\)，以及对每个 \(n\geqslant0\)，一个映射 \(\varphi_{n}\)（它从集合 \(C_{n}=J_{0}\times\cdots\times J_{n}\) 映到 X 的直径 \(\leqslant2^{-n}\) 的非空开闭子集组成的集合中），使得：

\begin{enumerate}
\def\labelenumi{(\roman{enumi})}
\item
  对每个 \(n \geqslant 0\)，存在整数 \(m_n \geqslant 1\)，使得 \(\mathrm{Card}(\mathrm{J}_n) = 2^{m_n}\)；
\item
  族 \((\varphi_0(c))_{c\in \mathrm{C}_0}\) 是 \(X\) 的一个分割；
\item
  对每个 \(n \geqslant 0\) 及每个 \(c \in \mathrm{C}_n\)，族 \((\varphi_{n+1}(c,j))_{j \in \mathrm{J}_{n+1}}\) 是 \(\varphi_n(c)\) 的一个分割。
\end{enumerate}

用 \(p_n\) 表示从 \(C_{n+1}\) 到 \(C_n\) 上的典范投影。序列 \(C = (C_n, p_n)_{n \geqslant 0}\) 是一个筛（TG, IX, p. 63, 定义 8）。与这个筛相联系的拓扑空间（TG, IX, p. 63）等同于拓扑空间 J，即诸离散拓扑空间 \(J_n\) 的乘积。筛 C 与映射序列 \((\varphi_n)_{n \geqslant 0}\) 定义了度量空间 X 的一个严格筛分（TG, IX, p. 63 与 p. 64）。由这个筛分得到的映射 \(f: J \to X\) 是连续的双射（TG, IX, p. 65）。由于拓扑空间 J 是紧的（TG, I, p. 63, 定理 3）而 X 是分离的，\(f\) 是同胚（TG, I, p. 63, 推论 2）。由于 \(J_n\) 同胚于 \(\{0, 1\}^{m_n}\)（\(n \geqslant 0\)），J 同胚于 \(\{0, 1\}^N\)（TG, I, p. 25, 命题 2）。

例. --- 设 K 为康托尔三分集（TG, IV, p. 9, 例）。对每个 \(n \geqslant 0\)，令 \(J_{n} = \{0,1\}\)，并如下定义映射 \(K_{n}\)（从集合 \(C_{n} = J_{0} \times \cdots \times J_{n}\) 到 {[}0,1{]} 的闭区间组成的集合中）：令 \(K_{0}(0) = [0, \frac{1}{3}]\) 与 \(K_{0}(1) = [\frac{2}{3}, 1]\)；对每个 \(n \geqslant 0\) 与每个 \(c \in C_{n}\)，\(K_{n+1}(c,0)\) 与

\(K_{n+1}(c,1)\) 分别是 \(K_n(c)\) 的“左三分”与“右三分”。若 \(c = (j_0,j_1,\ldots,j_n) \in C_n\)，则 \(K_n(c)\) 是（同前所引）中记作 \(K_{n,p}\) 的区间，其中 \(p = 2^n j_0 + 2^{n-1}j_1 + \cdots + j_n + 1\)；它也是区间 \([a,a + \frac{1}{3^{n+1}}]\)，其中 \(a = 2(\frac{j_0}{3} + \frac{j_1}{3^2} + \cdots + \frac{j_n}{3^{n+1}})\)。对 \(n \geqslant 0\) 与 \(c \in C_n\)，令 \(\varphi_n(c) = K_n(c) \cap K\)。族 \((\varphi_n(c))_{c \in C_n}\) 是由闭集组成的 K 的一个分割。因此这些集合在 K 中也是开的；它们非空且直径为 \(\frac{1}{3^{n+1}}\)，因为区间 \(K_n(c)\) 的端点都属于 K。序列 \(C = (C_n, p_n)_{n \geqslant 0}\)（其中 \(p_n: C_{n+1} \to C_n\) 是典范投影 \((c,j) \mapsto c\)）是一个筛，与这个筛相联系的拓扑空间等同于 \(\{0,1\}^N\)。筛 C 与映射序列 \((\varphi_n)_{n \geqslant 0}\) 定义了度量空间 K 的一个严格筛分。由这个筛分得到的映射 \(f: \{0,1\}^N \to K\) 是同胚，由下列公式给出

\[
f ((j _ {n}) _ {n \geqslant 0}) = 2 \sum_ {n = 0} ^ {\infty} \frac {j _ {n}}{3 ^ {n + 1}}.
\]

推论. --- 设 X 为可度量化、紧且非空的拓扑空间。存在从 \(\{0,1\}^{N}\) 到 X 上的连续满射。

每个紧可度量化的拓扑空间都同胚于拓扑空间 \(I^{N}\) 的一个子空间，且该子空间必然是闭的（TG, IX, p. 18, 命题 12 与 p. 21, 命题 16）。设 A 为 \(I^{N}\) 的闭子空间，h 为从 A 到 X 上的同胚。

令 \(K = \{0,1\}^{N}\)，且对 \(\alpha = (a_n) \in K\)，令 \(f(\alpha) = \sum_{n=0}^{\infty} a_n 2^{-n-1}\)；我们有 \(f(\alpha) \in I\)。这样定义的映射 \(f: K \to I\) 是满射（TG, IV, p. 42）且连续。事实上，若 K 的两个元素 \(\alpha\) 与 \(\beta\) 的指标 \(< n\) 的坐标都相同，我们有 \(|f(\beta) - f(\alpha)| \leqslant 2^{-n}\)。用 \(g\) 表示映射 \((\alpha_n) \mapsto (f(\alpha_n))\)（从 \(K^N\) 到 \(I^N\) 中）；它是连续的满射。拓扑空间 \(K^N\) 是紧的（TG, I, p. 63, 定理 3）、可度量化的（TG, IX, p. 15, 推论 2）且全不连通的（TG, I, p. 84, 命题 10），它的闭子空间 \(\bar{g}^{1}(A)\) 亦然；后者是非空的，因为映射 \(g\) 与 \(h\) 都是满射。

于是，空间 \(\bar{g}^{1}(A) \times K\) 同胚于 \(\{0,1\}^{N}\)（命题 14），因为它是紧、可度量化、全不连通且无孤立点的拓扑空间，而映射 \((x,y) \mapsto h(g(x))\) 从 \(\bar{g}^{1}(A) \times K\) 到 X 中是连续且满射的。

\subsection{道路像的拓扑性质}

命题 15. --- 分离拓扑空间中一条道路的像，是紧、可度量化、连通且局部道路连通的拓扑空间。

设 X 为分离拓扑空间，\(c\colon I \to X\) 为连续映射。用 R 表示 I 中的等价关系 \(c(s) = c(t)\)。拓扑空间 \(c(\mathbf{I})\) 是分离的（TG, I, p. 63, 推论 1），空间 I/R 是拟紧的（TG, I, p. 62, 定理 2），因此由 c 得到的双射 I/R \(\to c(\mathbf{I})\) 是同胚（TG, I, p. 63, 推论 2）。于是空间 \(c(\mathbf{I})\) 是紧的、可度量化的（TG, IX, p. 22, 命题 17）、连通的（TG, I, p. 82, 命题 6）且局部道路连通的（III, p. 261, 命题 8）。

定理 1（Hahn 与 Mazurkiewicz）. --- 每个可度量化、紧、非空、连通且局部连通的拓扑空间，都同胚于区间 {[}0,1{]} 的一个商空间。

引理 4. --- 设 K 为紧拓扑空间，R 为由覆盖 K 的 K 的开集组成的集合。存在 K 的一致结构的一个邻偶 V（TG, II, p. 27, 定理 1），使得对每个 \(x \in K\)，\(V(x)\) 都包含于 R 的某个元素之中。

对 K 的每个点 x，存在 K 的一致结构的一个邻偶 \(W_{x}\)，使得 \(\mathrm{W}_{x}(x)\) 包含于 R 的某个元素之中。设 \(V_{x}\) 为 K 的一致结构的满足下列条件的邻偶：\(\overset{2}{\mathrm{V}}_{x}\) 包含于 \(W_{x}\)。诸 \(\mathrm{V}_{x}(x)\) 的内部覆盖 K；由于空间 K 紧，存在 K 的有限子集 F，使得族 \((\mathrm{V}_{y}(y))_{y\in\mathrm{F}}\) 覆盖 K（TG, I, p. 59）。用 V 表示族 \((\mathrm{V}_{y})_{y\in\mathrm{F}}\) 的交；集合 V 是 K 的一致结构的邻偶。对每个点 \(x\in K\)，存在点 \(y\in F\)，使得 x 属于 \(\mathrm{V}_{y}(y)\)。于是集合 \(\mathrm{V}(x)\) 包含于 \(\overset{2}{\mathrm{V}}_{y}(y)\)，从而包含于 R 的某个元素之中。

引理 5. --- 设 X 与 Y 为一致空间，f 为从 X 到 Y 中的映射。设 F 为由 X 的闭子集组成的、覆盖 X 且具有下列性质的集族：

\begin{enumerate}
\def\labelenumi{(\roman{enumi})}
\item
  存在一个集合 \(F_{0} \in F\)，它与所有集合 \(F \in F\) 相交；
\item
  对 X 的一致结构的每个邻偶 U，不是 U 阶小集的集合 \(F \in F\) 只有有限多个；
\item
  对 Y 的一致结构的每个邻偶 V，集合 \(F \in F\)（使得 \(f(F)\) 不是 V 阶小集的）只有有限多个。
\end{enumerate}

那么，若 f 在每个集合 \(F \in F\) 上的限制都连续，则 f 连续。

设 x 为 X 的一点。我们来证明 f 在 x 处连续。存在集合 \(F_{1} \in F\)，使得 \(x \in F_{1}\)。\(f\) 在 \(F_{0} \cup F_{1}\) 上的限制是连续的（TG, I, p. 19, 命题 4）。把 \(F_{0}\) 换成 \(F_{0} \cup F_{1}\)，可化归为 \(x \in F_{0}\) 的情形。

设 V 为 Y 的一致结构的一个邻偶。选取该一致结构的一个邻偶 \(V'\)，使得 \(V' \subset V\)。由于 f 在 \(F_0\) 上的限制连续，存在 X 的一致结构的一个邻偶 U，使得 \(f(z) \in \mathrm{V}'(f(x))\)（对所有 \(z \in \mathrm{F}_0 \cap \mathrm{U}(x)\)）。设 \(U'\) 为 X 的一致结构的满足 \(U' \subset U\) 的一个邻偶。用 A 表示 \(F_0\)、诸集合 \(F \in F\) 中不是 \(U'\) 阶小集者、以及诸使 \(f(\mathrm{F})\) 不是 \(V'\) 阶小集的集合的并。按假设，A 是属于 F 的有限多个集合的并，而 f 在 A 上的限制连续（同前所引）。因此存在一个邻域 W（x 的，在 X 中），它包含于 \(U'(x)\)，且 \(f(y) \in \mathrm{V}(f(x))\)（对 \(y \in A \cap W\)）。为完成证明，只需证明也有 \(f(y) \in \mathrm{V}(f(x))\)（对每个点 \(y \in (\mathrm{X} - \mathrm{A}) \cap \mathrm{W}\)）。设 y 为这样一个点。设 F 为 F 的满足 \(y \in F\) 的一个元素。由 A 的定义，F 是 \(U'\) 阶小集，且 \(f(\mathrm{F})\) 是 \(V'\) 阶小集。按假设，F 与 \(F_0\) 相交。设 \(z \in F \cap F_0\)。我们有 \(z \in U'(y)\)，因为 F 是 \(U'\) 阶小集；又有 \(y \in U'(x)\)，因为 W 包含于 \(U'(x)\)；于是 \(z \in U'(x)\)，更有 \(z \in U(x)\)。但这时，由于 z 属于 \(F_0\)，我们有 \(f(z) \in V'(f(x))\)。此外 \(f(\mathrm{F})\) 是 \(V'\) 阶小集，于是 \(f(y) \in V'(f(z))\)。由此得出 \(f(y) \in V'(f(x))\)，最终有 \(f(y) \in V(f(x))\)。引理 5 证毕。

现在来证明定理 1。设 X 为非空的紧、连通且局部连通的度量空间。这样的空间是道路连通且局部道路连通的（III, p. 267, 推论 2）。由 III, p. 270 的推论与例，存在从康托尔三分集 K（TG, IV, p. 9, 例）到 X 上的连续满射 f。我们将构造一个连续开拓 g（f 的，到 {[}0,1{]} 上），这就证明了定理 1。

设 \(\varepsilon\) 为实数 \(>0\)。X 的直径 \(\leqslant \varepsilon\) 的开的道路连通子集覆盖 X。用 \(\mathcal{R}\) 表示它们在 \(f\) 下的原像组成的集合；这是由覆盖 K 的 K 的开子集组成的集合。由 III, p. 272 的引理 4，存在实数 \(\alpha > 0\)，使得 K 中半径为 \(\alpha\) 的每个闭球都包含于 \(\mathcal{R}\) 的一个元素之中。特别地，若 \(t\) 与 \(t'\) 为 K 中满足 \(|t - t'| \leqslant \alpha\) 的点，则存在 X 中连接 \(f(t)\) 与 \(f(t')\) 的道路，其像的直径 \(\leqslant \varepsilon\)。

前一段使我们可以归纳地构造严格递增序列 \((n_{k})_{k\geqslant0}\)，其项为整数 \(\geqslant0\)，且具有下列性质：对每个整数 \(k\geqslant0\) 与每对点 \((t,t')\)（K 的、满足 \(|t-t'| \leqslant3^{-n_{k}}\) 的），存在 X 中连接 \(f(t)\) 与 \(f(t')\) 的道路，其像的直径 \(\leqslant2^{-k}\)。K 在 {[}0,1{]} 中的余集是两两不交的开区间组成的族 \((\mathrm{I}_{n,p})\) 的并，其中 n 取遍整数 \(\geqslant0\) 组成的集合，而 p 取遍 1 与 \(2^{n}\) 之间的整数组成的集合（TG, IV, p. 9, 例）。考虑这些区间 \(I_{n,p}\) 之一，并把它写成 {]}a,b{[}。点 a 与 b 属于 K，且有 \(b-a=3^{-n-1}\)。如下定义区间 \(I_{n,p}\) 上的函数 g：选取 X 中连接 \(f(a)\) 与 \(f(b)\) 的道路 c，其像的直径 \(\leqslant2^{-k}\)（当 \(n_{k}\leqslant n<n_{k+1}\) 时），并令 \(g(t)=c(\frac{t-a}{b-a})\)（\(t\in I_{n,p}\)）。这样定义的函数 \(g\colon[0,1]\to X\) 是 f 的开拓。它在 K 上以及在每个闭区间 \(\overline{I_{n,p}}\) 上都连续。这些闭区间与 K 相交。此外，对每个实数 \(\varepsilon>0\)，区间 \(\overline{I_{n,p}}\) 中长度 \(>\varepsilon\) 的只有有限多个，而区间 \(\overline{I_{n,p}}\) 中在 g 下的像的直径 \(>\varepsilon\) 的也只有有限多个。由引理 5，映射 g 连续。
